Figs.~\ref{fig:haar_vs_sp} show that, in our current pipeline, the S+P transform does not surpass the Haar wavelet in rate--distortion performance. S+P remains competitive with Haar at very low bit-rates, but its PSNR degrades more quickly as quantization becomes coarser across all datasets. This behavior aligns with several structural differences between the two transforms and with the way our quantizer interacts with their coefficient statistics.

First, S+P modifies the basic Haar lifting step by introducing an in-loop prediction on the highpass branch. After computing the integer lowpass and preliminary detail coefficients $(s, d_1)$, S+P replaces $d_1[l]$ with a residual
\[
d[l] = d_1[l] - 
\left\lfloor 
    \frac{\mathrm{pred}(s, d_1[l+1])}{2^{\text{shift}}}
\right\rfloor,
\]
where $\mathrm{pred}$ is a fixed linear combination of neighboring lowpass differences and the next detail value. This predictor effectively acts as an additional highpass filter: it substantially reduces low-frequency leakage in the detail band whenever the correlation condition $\rho > 1/3$ holds, thereby lowering the variance and first-order entropy of $d[n]$. While this whitening is advantageous at the transform stage, it also concentrates more energy into the highest spatial frequencies, which are precisely the coefficients most heavily attenuated by our coarse scalar quantizer. Moderate step sizes that Haar tolerates well therefore produce noticeably larger PSNR losses under S+P.

Second, unlike Haar, which produces subbands with simple and predictable variance structure, the S+P transform applies additional prediction on the highpass branch. This prediction further suppresses low-frequency leakage in the detail bands, effectively shifting more residual energy toward the highest frequencies. Although this reduces the first-order entropy of the detail coefficients—an advantage for lossless or near-lossless coding—it also increases the sensitivity of these coefficients to coarse scalar quantization, which disproportionately attenuates high-frequency components. Because our quantization parameters were tuned heuristically and not optimized for S+P’s predictor-adjusted subbands, the resulting step sizes do not match the actual coefficient distributions well, contributing to the larger loss in PSNR at matched bit-rates.

Finally, Fig.~\ref{fig:effectEntropyCoding} shows that although S+P achieves some of the highest direct compression ratios before entropy coding, the additional gain from our backend is smaller for S+P than for Haar. Our entropy coder consists of a simple JPEG-style differential predictor followed by symbolwise Huffman coding; it does not perform run-length coding or adaptive context modeling. Because S+P already removes much of the spatial correlation and flattens the coefficient structure, it leaves less residual redundancy for this backend to exploit. Haar, by comparison, retains more correlation patterns that Huffman coding can capitalize on, leading to larger post–entropy-coding improvements.

Taken together, these factors explain why S+P appears promising at the transform stage, with lower entropy and strong direct compression, yet yields a weaker end-to-end CR--PSNR tradeoff than Haar under our current JPEG-style pipeline. With predictor-aware quantization and a more sophisticated entropy model, S+P’s relative position could plausibly change.
